\documentclass[../../main.tex]{subfiles}
\begin{document}

\begin{center}
  {\fontsize{48pt}{58pt}\selectfont \textbf{2001}\normalsize\textbf{年度}}
\end{center}

\begin{table}[htb]
  \centering
  \caption{出題テーマと難易度}
  \begin{tabular}{|c|p{9cm}|c|}
    \hline
    \textbf{問題番号} & \textbf{テーマ} & \textbf{難易度} \\
    \hline
    第1問 & 絶対値付き定積分の最小値と極限 & \\
    \hline
    第2問 & 空間図形（速さの異なる領域の到達範囲の体積） & \\
    \hline
    第3問 & 確率（カードの和がちょうど$k$になる確率） & \\
    \hline
    第4問 & 正方形の折り紙と線対称五角形の面積の最小値 & \\
    \hline
  \end{tabular}
\end{table}

\begin{table}[htb]
  \centering
  \caption{解答}
  \renewcommand{\arraystretch}{1.4}
  \begin{tabular}{|c|c|p{9cm}|}
    \hline
    \textbf{問題番号} & \textbf{小問} & \textbf{答え} \\
    \hline
    第1問 & (1) & $m(a)=\left(1-e^{-a/2}\right)^2$ \\
    \cline{2-3}
    & (2) & $\dfrac{1}{4}$ \\
    \hline
    第2問 & - & $V=\dfrac{\pi}{3}\left(2a^3+2+\dfrac{(a^2-1)\sqrt{a^2-1}}{a}\right)$ \\
    \hline
    第3問 & (1) & $P_N(1)=\dfrac{1}{N}$，$P_N(2)=\dfrac{1}{N}+\dfrac{1}{N^2}$，$P_N(3)=\dfrac{1}{N}+\dfrac{2}{N^2}+\dfrac{1}{N^3}$ \\
    \cline{2-3}
    & (2) & $P_3(4)=\dfrac{37}{81}$，$P_3(5)=\dfrac{121}{243}$ \\
    \cline{2-3}
    & (3) & $P_N(k)=\dfrac{1}{N}\left(1+\dfrac{1}{N}\right)^{k-1}$ \\
    \hline
    第4問 & - & $\sqrt{2}-1$ \\
    \hline
  \end{tabular}
\end{table}

\paragraph{試験形態}
4問．試験時間は要確認．

\end{document}
